\documentclass{article}
\usepackage{decomp}
\begin{document}
\decompnote{Collection markers: a cited work that is in the local collection is tagged \localcopy{<key>} after its reference entry; it is decompiled in decompiled/<key>/. A reference to another version of a local work (e.g. its preprint) is tagged \localother{<key>}. A reference without that tag is not in the local collection and needs an external pull. Tags added by \texttt{\detokenize{tools/mark_local.py}}.}

\sourcefile{1-s2.0-S0022227520369868-main.pdf}
\decompnote{Scanned PDF (ScienceDirect open archive, 4 pages). Pages 1, 2 and 4 are 1000 dpi bilevel scans and p.\,3 (with the Fig.~1 photograph) is a 300 dpi greyscale JPEG; the text layer is OCR. The page images were used as the authoritative text; OCR errors are corrected without notice here and listed in the README. The upper half of p.\,1 is blank in the scan (the article starts mid-page in the journal). Left out: the page footers ``Journal of Lipid Research Volume 17, 1976 Notes on Methodology'' with the page numbers 281--284; the footer is given once here.}

\pdfpage{1}
\noindent\textbf{Journal of Lipid Research}\quad Volume 17, 1976\quad \emph{Notes on Methodology}

\section*{Isolation of disaturated phosphatidylcholine with osmium tetroxide}

\noindent Robert J. Mason\textsuperscript{1}, Jean Nellenbogen, and John A. Clements\textsuperscript{2}

\noindent\emph{Cardiovascular Research Institute, University of California Medical Center, San Francisco, California 94143}

{\small\noindent\textsuperscript{1} Established Investigator of the American Heart Association.\\
\textsuperscript{2} Career Investigator of the American Heart Association.\par}
\decompnote{These two footnotes are printed at the foot of the left column of p.\,1; they are placed here after the author list.}

\begin{quote}\small
\textbf{Summary}\quad A simple, rapid, inexpensive method for isolating disaturated phosphatidylcholine from adult rat lung has been devised. Total lipids are reacted with osmium tetroxide dissolved in carbon tetrachloride, and the disaturated phosphatidylcholine is isolated on a column of neutral alumina. More than 99\% of the fatty acids in the phosphatidylcholine fraction are saturated and 94\% of this material migrates as phosphatidylcholine on subsequent thin-layer chromatography.

\textbf{Supplementary key words}\quad lung lipids · phosphatidylcholine
\end{quote}

For investigators interested in phospholipid metabolism in the lung, there is a need for a rapid method for isolating disaturated phosphatidylcholine. This species of phosphatidylcholine, which has only saturated fatty acids, is found in high concentration in purified surface active material and is mainly responsible for the ability of the latter substance to produce low surface tensions (1–3). Ever since dipalmitoylphosphatidylcholine was first isolated from lung by Thannhauser, Benotti, and Boncoddo (4), many means of isolating disaturated phosphatidylcholine have been devised. These methods include precipitation in cold ethanol (1) or acetone (4, 5), oxidation with ozone (6) or permanganate and periodate (7), formation of mercuric acetate adducts (8–10), and argentation thin-layer chromatography (11). Species of phosphatidylcholines can also be separated by diglycerides after treatment with phospholipase C (12, 13). Most of these methods require isolation of phosphatidylcholines before the specific separation procedures. Exceptions are the method of Young and Tierney (14), in which mercuric acetate adducts are formed with total phospholipids and the disaturated phosphatidylcholine is isolated by thin-layer chromatography, and the method of Brown (1) that uses solvent fractionation for analysis of material obtained by endobronchial lavage.

A rapid, simple, quantitative method for isolating disaturated phosphatidylcholine would be useful for analyzing amniotic fluid to assess fetal maturity, for studying secretion of surface active material by type II alveolar cells, and for performing turnover studies in intact lung. We have developed a method which does not require prior isolation of phosphatidylcholine from lipid extracts.

Osmium tetroxide is a powerful oxidant that reacts with double bonds to form complexes (15–18). One molecule of osmium tetroxide reacts with each double bond in egg phosphatidylcholine, and the reaction proceeds rapidly to completion (15, 17).

Our method is simply to react total lipids with osmium tetroxide in carbon tetrachloride and to isolate the disaturated phosphatidylcholine by column chromatography on neutral alumina.

\subsection*{Methods}

Lipids are prepared by extracting cells or homogenized tissue with chloroform–methanol 2:1 (v/v), \pdfpage{2}filtering, separating the solvents into two phases by adding 0.1 M KCl, and washing according to the procedure of Folch, Lees, and Sloane Stanley (19). Osmium tetroxide (Electron Microscopy Sciences, Port Washington, Pa.) is dissolved in carbon tetrachloride (100 mg/ml) and stored in the dark at 4°C. This stock solution is stable for at least two years. An aliquot of lipid (1 mg) is evaporated and the residue is redissolved in 0.5 ml of carbon tetrachloride that contains 3.1 mg of osmium tetroxide. The reacting solution darkens, presumably as osmium dioxide is formed (15). After waiting 15 min for completion of the reaction, the solution is evaporated and the residue is redissolved in chloroform–methanol 20:1 (v/v). This material is applied to a column formed by placing 0.8 g of aluminum oxide (neutral alumina, 100–200 mesh, Biorad Laboratories, Richmond, California, lots \#10461 and \#7579, activated for 1 hr at 110°C) on a plug of glass wool in the neck of a disposable 9-in Pasteur pipette. Neutral alumina was chosen as the adsorbent for chromatography because it binds acidic lipids under the conditions we use in this method. Neutral lipids not altered by the osmium tetroxide are first eluted with 10 ml of chloroform–methanol 20:1 (v/v). The disaturated phosphatidylcholine is then eluted with 5 ml of chloroform–methanol–7 M ammonium hydroxide 70:30:2 (v/v). A dark residue remains at the top of the column.

\decompnote{Table 1 is printed at the top of the left column of p.\,2, before the text that continues from p.\,1; it is placed here, at the end of that first paragraph.}
\begin{center}
\noindent TABLE 1. Recovery of disaturated phosphatidylcholine\par\smallskip
\begin{tabular}{llrcc}
\toprule
Method & Preparation & N & Disaturated PC\textsuperscript{\emph{b}} \% of Total PL & Recovery (\%) \\
\midrule
Osmium tetroxide & A\textsuperscript{\emph{a}} & 10 & 23.1 ± 0.4 & 90.7 ± 2.8 \\
Osmium tetroxide & B\textsuperscript{\emph{a}} & 6 & 21.6 ± 0.4 & 92.5 ± 1.2 \\
Mercuric acetate adducts & A & 3 & 21.6 ± 1.8 & 87.3 ± 1.0 \\
\bottomrule
\end{tabular}
\end{center}
{\small\noindent\textsuperscript{\emph{a}} Lipids extracted from lungs of two different rats (A and B).\\
\textsuperscript{\emph{b}} The disaturated phosphatidylcholine (PC) is expressed as a percentage of total phospholipid (PL) as determined by measurement of total lipid phosphorus (20). The recoveries are determined from the added \textsuperscript{14}C-labeled dipalmitoylphosphatidylcholine (21). The numbers are the mean ± standard deviation.\par}
\decompnote{The table prints ``Preparation'' hyphenated over two lines (``Prepara-/tion'') and the method names over two lines (``Osmium/tetroxide'', ``Mercuric acetate/adducts''). Also in tables/table1.csv.}

The reaction of unsaturated lipid and osmium tetroxide proceeds rapidly and is complete in the time necessary to evaporate the sample, but we have arbitrarily decided to wait 15 min for the reaction to go to completion. When the columns are eluted, the black band should stay at the top and will do so unless the rate of elution is too fast. The rate of elution (approximately 0.8 ml/min) is controlled by the packing of the glass wool in the neck of the Pasteur pipette. This detail is easily learned by experience. Up to 12 columns have been run simultaneously and they can stand for as long as 15 min between eluting solutions without affecting the separation or recovery.

The methods used in the analyses have been reported previously (10, 20–22). Phosphorus was determined by the method of Bartlett (20). Radioactivity was measured by scintillation spectroscopy with samples prepared with Aquasol (New England Nuclear Corp., Boston, Mass.)—water 15:4 (v/v). Methyl esters of fatty acids of total lipids and isolated phosphatidylcholines were prepared with methanol and sulfuric acid, separated by gas–liquid chromatography on a column of ethylene glycol succinate, and quantitated by integration of the areas of the individual peaks (21). The purity of the fraction that contained the isolated disaturated phosphatidylcholine was determined by thin-layer chromatography and phosphorus measured in the sample lanes and in blank lanes (22).

Isolation of disaturated phosphatidylcholine by the mercuric acetate adduct method, for comparison with the method in this report, required the isolation of phospholipids by silicic acid column chromatography, isolation of mixed phosphatidylcholines by thin-layer chromatography, and isolation of disaturated phosphatidylcholine by silicic acid column chromatography after formation of mercuric acetate adducts of the unsaturated species (21).

\subsection*{Results and discussion}

We used total lipids extracted from rat lung to develop the method. The results are shown in \textbf{Tables 1} and \textbf{2}. \textsuperscript{14}C-labeled dipalmitoylphosphatidylcholine (synthesized chemically in our laboratory) was added to the total lipids for quantitation of recovery (21). Two batches of rat lung lipids were tested and gave similar results. The values agreed with those obtained by a mercuric acetate adduct method (21). The recoveries were 90\% with 16 samples of rat lung and 90.3 ± 3.1\% (mean ± standard deviation) with a total of 71 samples of various tissue extracts. All the fatty acids (>99\%) from phosphatidylcholines isolated from rat and rabbit lungs by this method have been saturated (N = 30). As shown in \textbf{Fig. 1}, the material isolated by this method appears as a single spot in the phosphatidylcholine area on thin-layer chromatography, and this spot contains 93.9 ± 1.5\% (N = 3) of the phosphorus recovered from the plate (22). The limits of the resolution of the thin-layer chromatographic separation is 98\%, as judged by the distribution of purified, chemically synthesized [\textsuperscript{14}C]dipalmitoylphosphatidylcholine. Other disaturated phospholipids could be present in trace amounts.

\pdfpage{3}
\begin{center}
\noindent TABLE 2.\quad Fatty acid composition\textsuperscript{\emph{a}}\par\smallskip
\small
\begin{tabular}{lccccccccc}
\toprule
Preparation & 14:0\textsuperscript{\emph{b}} & 15:0 & 16:0 & 16:1 & 18:0 & 18:1 & 18:2 & 20:4 & Unidentified \\
\midrule
A & 1.9 & tr\textsuperscript{\emph{c}} & 35.8 & 4.6 & 11.8 & 14.6 & 18.5 & 10.4 & 2.5 \\
B & 2.5 & tr & 39.5 & 5.2 & 10.8 & 12.4 & 10.8 & 14.4 & 4.1 \\
Disaturated phosphatidylcholine & 6.3 & 1.9 & 87.9 & & 3.9 & & & & \\
\bottomrule
\end{tabular}
\end{center}
{\small\noindent\textsuperscript{\emph{a}} The mean of triplicate analyses of methyl esters of total lipids and the mean of 16 determinations of the isolated disaturated phosphatidylcholines are presented. There was no difference in the fatty acid composition of the disaturated phosphatidylcholine of the two preparations (A and B). The percent composition is given as mole \%.\\
\textsuperscript{\emph{b}} Number of carbon atoms: number of double bonds.\\
\textsuperscript{\emph{c}} tr stands for trace (0.5\%).\par}
\decompnote{Table 2 spans both columns at the top of p.\,3. ``Disaturated phosphatidylcholine'' is printed over two lines; its cells for 16:1, 18:1, 18:2, 20:4 and Unidentified are blank in the print. Also in tables/table2.csv.}

Under different conditions of chromatography and with different solvents, acidic phospholipids can be eluted from columns of neutral alumina and some phosphatidylcholine can be converted to lysophosphatidylcholine (23, 24). With our procedure and lipids from adult rat lung, no acidic lipids or lysolecithin were eluted as determined by inspection of the thin-layer plate (Fig. 1) or measuring the phosphorus on the plate.

We have shown that for adult rat lung 94\% of the phosphorus found in the chloroform–methanol–ammonium hydroxide eluate is disaturated phosphatidylcholine. Similar results have been found with fluid obtained by endobronchial lavage of lungs of adult rabbits.\textsuperscript{3} For use with other biologic materials, the method will have to be validated. We have not tested this method with amniotic fluid samples. Our method is less satisfactory with extracts that contain very little disaturated phosphatidylcholine. With rabbit liver and heart, 3.9\% and 4.2\% of the total phospholipid elutes as disaturated phosphatidylcholine. It is likely that some of the phosphorus eluted with chloroform–methanol–ammonium hydroxide is not disaturated phosphatidylcholine. Previously, with different samples of these tissues, we found that ≤2\% of the phospholipids were disaturated phosphatidylcholine. The method used required initial isolation of mixed phosphatidylcholines and then separation of the disaturated species after mercuric acetate adduction (25). For samples containing a low percentage of disaturated phosphatidylcholine (<10\% of lipid phosphorus), this species can be separated from the unsaturated phosphatidylcholines by thin-layer chromatography on silica H plates with chloroform–methanol–7 M ammonium hydroxide 70:30:3.5 (v/v) after reaction of isolated phosphatidylcholines (but not total lipids) with osmium tetroxide.

The method as described in this report will be valuable for multiple analyses of similar samples for comparative purposes (i.e. secretion of surface active material) and for simple analysis of lipids containing a moderate amount of disaturated phosphatidylcholine (>10\% of lipid phosphorus). The composition of the eluate from the alumina column should be checked with different biologic materials, and the product may have to be further purified for materials that contain a low amount of disaturated phosphatidylcholine.

{\small\noindent\textsuperscript{3} Oyarzun, M., and J. A. Clements, unpublished observations.\par}
\decompnote{Footnote 3 is printed at the foot of the left column of p.\,3.}

\begin{center}
\includegraphics[width=0.6\linewidth,height=0.55\textheight,keepaspectratio]{figures/fig1.png}
\end{center}
\noindent\textbf{Fig. 1.}\quad A thin-layer plate of lipid standards, rat lung lipids, and eluates from the alumina column. Different lipids were applied to the designated lanes: (A) 375 μg of a phospholipid standard (phosphatidylethanolamine, phosphatidylcholine, sphingomyelin, and lysophosphatidylcholine); (B) 340 μg of total lipids from rat lung; (C) 310 μg of rat lung phospholipids eluted with chloroform–methanol–ammonium hydroxide (not reacted with osmium tetroxide); (D) 200 μg of the same phospholipids after reaction with osmium tetroxide, (E) blank lane for background phosphorus measurement, (F) same as A. 94.9\% of the phosphorus in lane D is in the phosphatidylcholine area, 3.7\% in the sphingomyelin area, 0\% in the lysophosphatidylcholine area and the remainder at the origin or above the phosphatidylcholine area. The silica H thin-layer plate was developed with chloroform–methanol–acetic acid–water 100:50:14:7 (v/v), sprayed with 0.25\% 8-anilino-1-naphthalenesulfonic acid in water and illuminated with ultraviolet light.
\alttext{Fig.~1 is at the lower right of p.\,3. Black-and-white photograph of a thin-layer plate under ultraviolet light: a dark plate with bright (fluorescent) spots, divided by thin dark vertical scored lines into lanes; the six sample lanes are labelled A, B, C, D, E, F below the photograph (sans-serif capitals). The plate's origin and solvent front are not marked. Positions below are measured on the 300 ppi page image as the fraction of the photograph's height from its bottom edge (work/lanes.py; band extents where the brightness is above half-way between background and saturation; about ±0.02). Lane A (phospholipid standard): a faint, diffuse spot near the top (centre ≈0.84) and three bright, saturated bands low on the plate at ≈0.30–0.41, ≈0.16–0.27 and ≈0.08–0.14, separated by darker gaps. Lane B (total rat lung lipids): a faint spot near the top (≈0.83–0.90, at about the height of lane A's top spot), a moderately bright spot at ≈0.53–0.63, a large bright band at ≈0.31–0.47 and a bright band at ≈0.18–0.28; a dim smear runs between them. Lane C (eluted phospholipids, not reacted with osmium tetroxide): one large, very bright elongated region from ≈0.21 to ≈0.53 with a narrow darker waist at about 0.37, i.e.\ two touching spots, the upper one larger. Lane D (the same phospholipids after osmium tetroxide): a single bright spot at ≈0.32–0.47, at about the height of the upper part of lane C's region and of lane A's highest bright band, with a very faint tail below it; nothing else is visible in the lane. Lane E (blank): no spots. Lane F (same as A): the same pattern as lane A, a faint spot near the top (≈0.83) and three bright bands at ≈0.32–0.40, ≈0.19–0.27 and ≈0.10–0.15. Which band in lanes A and F is which standard is not labelled in the figure.}

\pdfpage{4}
During the development of this method, several other means of separation and refinements of the technique were evaluated. With an eluting solution of chloroform–methanol 1:1 (v/v), the recovery decreased to 75\%. Addition of acetic acid to the solution elutes excess phosphorus, some of the black material, and fatty acids. Separation on columns of diethylaminoethyl cellulose or thin-layer plates of silica gel or alumina were less satisfactory in terms of both purity of the final product and of total recovery.

This method is similar in principle to a recently published procedure by Shimojo, Abe, and Ohta (7) in which mixed phosphatidylcholines were oxidized with permanganate and periodate and the unaffected disaturated phosphatidylcholine was isolated by thin-layer chromatography. The major differences are that the permanganate and periodate method apparently requires an initial isolation of phosphatidylcholines from mixed lipids, an extraction procedure after the reaction, thin-layer chromatography for isolation of the disaturated phosphatidylcholine, and elution from the thin-layer plate for recovery of the disaturated phosphatidylcholine. Our method is quicker, simpler, and just as reliable for processing whole lung lipids.

In conclusion, a rapid method for isolating disaturated phosphatidylcholine from mixed lipids from adult rat lung has been devised. A drawback of the method is that unsaturated lipids are probably destroyed and cannot be analyzed further.
\decompnote{The paragraph ends with the journal's end-of-article mark, a small black square with the letters ``JLR''.}

{\small\noindent This work was supported in part by grants HL-06285 and HL-05705 from the National Heart and Lung Institute.

\noindent\emph{Manuscript received 18 August 1975; accepted 2 February 1976.}\par}

\begin{center}
REFERENCES
\end{center}

\begin{hangpars}\small
1. Brown, E. S. 1964. Isolation and assay of dipalmityl lecithin in lung extracts. \emph{Amer. J. Physiol.} \textbf{207:} 402–406.

2. King, R. J. 1974. The surfactant system of the lung. \emph{Federation Proc.} \textbf{33:} 2238–2247.

3. Goerke, J. 1974. Lung surfactant. \emph{Biochim. Biophys. Acta.} \textbf{344:} 241–261.

4. Thannhauser, S. J., J. Benotti, and N. F. Boncoddo. 1946. Isolation and properties of hydrolecithin (dipalmityl lecithin) from lung; its occurrence in the sphingomyelin fraction of animal tissues. \emph{J. Biol. Chem.} \textbf{166:} 669–675.

5. Gluck, L., E. K. Motoyama, H. L. Smits, and M. V. Kulovich. 1967. The biochemical development of surface activity in mammalian lung. I. The surface-active phospholipids; the separation and distribution of surface-active lecithin in lung of the developing rabbit fetus. \emph{Pediatr. Res.} \textbf{1:} 237–246.

6. Privett, O. S., and M. L. Blank. 1963. A method for the structural analysis of triglycerides and lecithins. \emph{J. Amer. Oil Chem. Soc.} \textbf{40:} 70–75.

7. Shimojo, T., M. Abe, and M. Ohta. 1974. A method for determination of saturated phosphatidylcholine. \emph{J. Lipid Res.} \textbf{15:} 525–527.

8. Blank, M. L., L. J. Nutter, and O. S. Privett. 1966. Determination of the structure of lecithins. \emph{Lipids.} \textbf{1:} 132–135.

9. King, R. J., and J. A. Clements. 1970. Separation by gel chromatography of naturally occurring phosphatidylcholine mixtures according to number of ethylenic linkages. \emph{J. Lipid Res.} \textbf{11:} 381–385.

10. Mason, R. J., G. Huber, and M. Vaughan. 1972. Synthesis of dipalmitoyl lecithin by alveolar macrophages. \emph{J. Clin. Invest.} \textbf{51:} 68–73.

11. Spitzer, H. L., J. R. Norman, and K. Morrison. 1969. In vivo studies of [Me-\textsuperscript{3}H]choline and [1, 2-\textsuperscript{14}C\textsubscript{2}]choline incorporation into lung and liver lecithins. \emph{Biochim. Biophys. Acta.} \textbf{176:} 584–590.

12. Kuksis, A., L. Marai, W. C. Breckenridge, D. A. Gornall, and O. Stachnyk. 1968. Molecular species of lecithins of some functionally distinct rat tissues. \emph{Can. J. Physiol. Pharmacol.} \textbf{46:} 511–524.

13. Tierney, D. F., J. A. Clements, and H. J. Trahan. 1967. Rates of replacement of lecithins and alveolar instability in rat lungs. \emph{Amer. J. Physiol.} \textbf{213:} 671–676.

14. Young, S. L., and D. F. Tierney. 1972. Dipalmitoyl lecithin secretion and metabolism by rat lung. \emph{Amer. J. Physiol.} \textbf{222:} 1539–1544.

15. Stoekenius, W., and S. C. Mahr. 1965. Studies on the reaction of osmium tetroxide with lipids and related compounds. \emph{Lab. Invest.} \textbf{14:} 458–469.

16. Korn, E. D. 1967. A chromatographic and spectrophotometric study of the products of the reaction of osmium tetroxide with unsaturated lipids. \emph{J. Cell Biol.} \textbf{34:} 627–638.\localcopy{korn1967}

17. Riemersma, J. C. 1968. Osmium tetroxide fixation of lipids for electron microscopy: a possible reaction mechanism. \emph{Biochim. Biophys. Acta.} \textbf{152:} 718–727.\localcopy{riemersma1968}

18. Collin, R., W. P. Griffith, F. L. Phillips, and A. C. Skapski. 1973. Staining \emph{and fixation of unsaturated} membrane lipids by osmium tetroxide; crystal structure of a model osmium (VI) intermediate. \emph{Biochim. Biophys. Acta.} \textbf{320:} 745–747.\localcopy{collin1973}

19. Folch, J., M. Lees, and G. H. Sloane Stanley. 1957. A simple method for the isolation and purification of total lipids from animal tissues. \emph{J. Biol. Chem.} \textbf{226:} 497–509.

20. Bartlett, G. R. 1959. Phosphorus assay in column chromatography. \emph{J. Biol. Chem.} \textbf{234:} 466–471.

21. King, R. J., and J. A. Clements. 1972. Surface active materials from dog lung II. Composition and physiological correlations. \emph{Amer. J. Physiol.} \textbf{223:} 715–726.

22. Parker, F., and R. F. Peterson. 1965. Quantitative analysis of phospholipid and phospholipid fatty acids from silica gel thin-layer chromatograms. \emph{J. Lipid Res.} \textbf{6:} 455–465.

23. Renkonen, O. 1962. Breakdown of lecithin on aluminum oxide columns. \emph{J. Lipid Res.} \textbf{3:} 181–183.

24. Luthra, M. G., and A. Sheltawy. 1972. The chromatographic separation of phospholipids on alumina with solvents containing ammonium salts. \emph{Biochem. J.} \textbf{126:} 251–253.

25. Mason, R. J. 1973. Disaturated lecithin concentration of rabbit tissues. \emph{Amer. Rev. Respir. Dis.} \textbf{107:} 678–679.
\end{hangpars}
\decompnote{Ref.~18: the words ``and fixation of unsaturated'' are printed in italic, as reproduced; the rest of the title is roman.}

\bigskip
\noindent This is an Open Access article under the CC BY license.
\decompnote{This line is printed (as real text, with ``CC BY'' as a link) below the footer of p.\,1 by the publisher; it is the PDF's only licence statement.}

\end{document}
